\documentclass[xcolor=dvipsnames]{beamer}
\usetheme{Rochester}
\usepackage{graphicx}
\usecolortheme[named=OliveGreen]{structure}
%\useoutertheme{infolines}
\useinnertheme{rounded}
\setbeamertemplate{blocks}[rounded][shadow=false]
\setbeamertemplate{items}[ball]
\begin{document}
\title[Netsoc Git Tutorial] {Dublin University Internet Society}
\author[A. Anderson \& C. Vize]{Andrew Anderson \inst{1} \and Colm Vize \inst{2}}
\institute[TCD]{\inst{1} School Of Computer Science and Statistics \and \inst{2} Tophat Software}
\subtitle{Git Tutorial}{An introduction to the git VCS}
\date{\today}
\titlepage

\begin{frame}
	\frametitle{A tutorial introduction}

	Lets start by creating a repo
	\begin{block}{Create a new repo}
		mkdir myrepo\\
		cd myrepo\\
		git init
	\end{block}
\end{frame}

\begin{frame}
	\frametitle{Cloning repositories}
	Makes a local copy of a repository, form e.g github\\
	Works with lot of protocols (http, ssh, native git protocol, even local files).

	\begin{block}{Clone a local repository}
		git clone /path/to/repo
	\end{block}
	\begin{block}{Clone a remote repository (These slides!)}
		git clone https://github.com/netsoc/Git-Tutorial.git
	\end{block}
\end{frame}

\begin{frame}
	\frametitle{Make your changes}
	After the cloning a repo, as on the prevois slides, you will have a folder with the repositories
	files in it (normally source code for a program).

	Make your changes, and then...
\end{frame}

\begin{frame}
	\frametitle{Add and commit}
	Commit them.
	You can stage chages (select them for committing) by using
	\begin{block}{}
		git add $\langle$filename$\rangle$ \# adds $\langle$filename$\rangle$\\
		git add . \# also works on directories
	\end{block}
	To actually commit these changes use
	\begin{block}{}
		git commit -m "Commit message"
	\end{block}
	commits the changes we have added, with the message "Commit Message"\\
	
	\begin{block}{}	
		git commit
	\end{block}
 	If you don't specify -m "whatever", git will open a text editor, and you can type your commit message in there 
	Now the file is committed to your local repo, but not on the remote repo \emph{yet}.
\end{frame}

\begin{frame}
	\frametitle{A tutorial introduction}

	\begin{block}{Workflow}
		Your local repository consists of three "trees" maintained by git. The first one is 
		your \emph{Working Directory} which holds the actual files. the second one is the \emph{Index} which acts as a staging area and finally the \emph{HEAD} which points to the last commit you've made.
	\end{block}
	\begin{center}
		\includegraphics[scale=0.3]{trees.png}
	\end{center}
\end{frame}


\end{document}
